\begin{abstract}
For every fixed treewidth bound, the shortest-path metrics of finite
connected graphs with arbitrary positive real edge lengths embed into
real $L_1$ with uniformly bounded distortion. This resolves the
bounded-treewidth case of the Gupta--Newman--Rabinovich--Sinclair
conjecture positively.
\end{abstract}
\tableofcontents
\clearpage

\section{Introduction}\label{sec:introduction}
A graph metric is the shortest-path metric $d_{G,\ell}$ of a finite
connected undirected graph $G=(V,E)$ with positive real edge lengths
$\ell:E\to(0,\infty)$. An embedding into real $L_1$ has
\emph{distortion at most $C$} if it can be rescaled so that every
image distance is between the original distance and $C$ times that
distance. We establish a bound depending only on the treewidth.

A \emph{tree decomposition} consists of a finite tree $T$ and subsets
$B_t\subseteq V$, called bags, indexed by $t\in V(T)$: the bags cover
$V$, each graph edge has both endpoints in a bag, and the bags
containing any fixed vertex induce a connected subtree of $T$.
Its width is $\max_t|B_t|-1$; the treewidth of $G$ is the minimum
width of a tree decomposition. Throughout this paper $k$ bounds
\emph{bag cardinality}, so the corresponding treewidth bound is $k-1$.

\begin{theorem}\label{thm:main}
For every integer $k\ge2$ there is a finite constant $C_k\ge1$ with the
following property. If a finite connected graph $G=(V,E)$ admits a tree
decomposition with bags of size at most $k$, then, for every edge-length
function $\ell:E\to(0,\infty)$, there is a map $F$ from $V$ into a real
$L_1$ space such that
\begin{equation}\label{eq:main-embedding}
 d_{G,\ell}(u,v)
 \le \|F(u)-F(v)\|_1
 \le C_k\,d_{G,\ell}(u,v)
 \qquad(u,v\in V).
\end{equation}
The constant is independent of the number of vertices, the depth of the
decomposition, and all ratios of edge lengths.
\end{theorem}

The target can be taken to be a finite-dimensional real $\ell_1$ space.
The proof gives the existence of $C_k$, without an explicit estimate
for its dependence on $k$. The bounded-treewidth family is one of the
cases of the conjecture of Gupta, Newman, Rabinovich, and Sinclair,
which predicts a uniform $L_1$ distortion bound for every proper
minor-closed graph family~\cite{GNRS}. By the excluded-planar-minor
theorem of Robertson and Seymour~\cite{RS86Planar}, our result also
gives a uniform bound for graphs excluding any fixed planar graph
as a minor.

\subsection{Context and prior work}

Finite $L_1$ pseudometrics are exactly the nonnegative combinations of cut
pseudometrics. If $S\subseteq V$ and
$\delta_S(u,v)=|\mathbf1_S(u)-\mathbf1_S(v)|$, then
\begin{equation}\label{eq:intro-cut-representation}
 D(u,v)=\sum_{S\subseteq V}w_S\delta_S(u,v),\qquad w_S\ge0,
\end{equation}
is realized by the vectors $(w_S\mathbf1_S(v))_{S\subseteq V}$.
Conversely, threshold sweeps through real coordinates give this
representation for any finite set of $L_1$ points. We therefore
seek cuts that charge graph edges by at most a constant times their
lengths and separate every vertex pair by a comparable amount.
The edge bound extends along shortest paths to all pairs.

Linial, London, and Rabinovich related low-distortion embeddings to
approximate multicommodity flow--cut bounds~\cite{LLR95}; Aumann and Rabani independently
obtained a general approximate max-flow/min-cut theorem~\cite{AR98}. Given nonnegative edge capacities and nonnegative demands on unordered pairs of distinct
vertices, not all zero, let $\lambda_*$ be the largest common multiplier of
the demands that can be routed fractionally, and let $\phi_*$ be the minimum
ratio of crossing capacity to separated demand over cuts with
positive separated demand. For instances with $\lambda_*>0$, the
ratio $\phi_*/\lambda_*$ is the flow--cut gap. Minimizing the cut
ratio that defines $\phi_*$ is the nonuniform sparsest-cut problem. The worst such gap
on a fixed graph equals the worst $L_1$ distortion of its weighted
shortest-path metrics~\cite[Section~3, Theorem~3.2]{GNRS}.
For graphs with bags of size at most $k$, Theorem~\ref{thm:main}
therefore gives
\[
 \lambda_*\leq\phi_*\leq C_k\lambda_*.
\]

Gupta et al. proved the treewidth-two case~\cite{GNRS}.
The optimal distortion is $2$, by the upper bound of Chakrabarti,
Jaffe, Lee, and Vincent~\cite{CJLV08} and the matching lower bound
of Lee and Raghavendra~\cite{LR10}. Other structural restrictions
also yielded uniform bounds: a fixed number of successive outer-face
removals in a planar graph~\cite{CGNRS06}, and graphs formed by
identifying edges of pieces from a fixed finite family~\cite{LP13}.
A path decomposition is a tree decomposition whose indexing tree is
a path; its minimum width is the graph's \emph{pathwidth}.
Lee and Sidiropoulos established the bounded-pathwidth case through
stochastic embeddings into trees: each sampled tree metric dominates
the graph metric, while its expected expansion is bounded by a function
of the pathwidth~\cite{LeeSidiropoulos}.
Abraham, Filtser, Gupta, and Neiman subsequently obtained
$O(\sqrt p)$ distortion into $L_1$ for pathwidth $p$~\cite{AFGN22}.
The route through dominating trees cannot yield a uniform bound for
all bounded-treewidth graphs: some series-parallel metrics require
expected tree distortion $\Omega(\log n)$~\cite[Theorem~5.6]{GNRS}.
We instead construct cut measures directly on a tree decomposition,
allowing arbitrary branching.

A closer precursor is the local-consistency theorem of Chlamt\'{a}\v{c},
Krauthgamer, and Raghavendra~\cite[Corollary~1.2]{CKR}. For treewidth
$t$, they obtain bounded $L_1$ distortion when the restrictions of the
metric to sets of $t+3$ vertices admit isometric cut representations
whose distributions agree on overlaps. Their work supplies both
the Markov gluing of consistent local laws and the bounded-state
Markov-flow theorem used below. Here we construct the bag laws
directly and prove separation without assuming such locally
consistent isometric representations of the original metric.

The same line of work developed approximation algorithms for
nonuniform sparsest cut using stronger relaxations~\cite{CKR,
GuptaTalwarWitmer,ChalermsookEtAl,CohenAddadMomkeVerdugo}.
These contain consistency information beyond the metric relaxation
whose gap is measured by $L_1$ distortion, so their guarantees do
not by themselves give the embedding theorem. There is nevertheless
an algorithmic consequence of our result. Gupta, Talwar, and Witmer
also approximate the minimum $L_1$ distortion~\cite[Section~3.3.1]{GuptaTalwarWitmer},
and Cohen-Addad, M\"omke, and Verdugo obtain fixed-parameter running
time for this task~\cite{CohenAddadMomkeVerdugo}. For positive rational
edge lengths, these algorithms and Theorem~\ref{thm:main} yield
embeddings of distortion at most $(2+\varepsilon)C_k$, for each fixed
$\varepsilon>0$, in fixed-parameter time with respect to $k$.

The other published input is the padded-decomposition theorem of
Filtser et al.~\cite{FiltserEtAl}. At each scale, it supplies a
random partition into small-diameter parts with positive probability
that a small ball around any specified vertex stays in one part.
The padding parameters depend only on the treewidth. The same work
gives the size-dependent $L_1$ distortion bound
$O(\sqrt{\log(2+t)\log n})$ for $n$-vertex graphs of treewidth $t$,
through its Euclidean embedding estimate~\cite{FiltserEtAl}. Such partitions
allow localized Lipschitz perturbations. Conroy and Filtser obtained related logarithmic
padding bounds for general excluded-minor classes~\cite{ConroyFiltser25}.
The issue addressed here is how to retain useful perturbations without accumulating a loss over
the number of scales or the depth of the decomposition.

The companion paper on planar graph metrics~\cite{OpenAIPlanar}
proves a universal distortion bound for that class. It also combines
its planar theorem with Theorem~\ref{thm:main}
and the stochastic reduction of Lee and Sidiropoulos~\cite[Theorem~1.7]{LS09Geometry}
to obtain uniform $L_1$ distortion and nonuniform fractional flow--cut
bounds for each fixed almost-embeddable family. Roughly, these graphs
consist of a part embedded in a fixed surface, boundedly many additional
pieces whose bounded-width path decompositions follow face-boundary
order, and boundedly many extra vertices with arbitrary adjacencies.
The reduction uses $1$-sums, or gluings at single vertices, of planar
and bounded-treewidth graphs. Extending this conclusion through
gluings along larger complete subgraphs remains a further step toward
the full conjecture.

\subsection{The proof: preserving local laws and transporting separation}

For a bag $B$, a \emph{bag cut law} assigns nonnegative masses to the
binary labelings $\{0,1\}^B$. Two neighboring laws are consistent
when their restrictions to the intersection agree. Because the bags
form a tree, consistent laws glue into a global law of cuts: after
sampling one bag, sample each neighboring bag conditionally on the
labels of its intersection. Section~\ref{sec:cuts} gives this
Markov gluing and its exact marginals. Controlling the separation
mass of an edge in its containing bag gives the global upper bound.
Consistency alone does not supply a lower bound for vertices in
distant bags.

\paragraph{Particles preserve bag laws under repeated updates.}
A \emph{particle} is a real function on all graph vertices. Sweeping
a threshold through its values gives a cut measure. We keep a fixed
number of particles and sum their symmetric threshold measures,
which give each cut and its complement equal mass.
Suppose two particles have the same weak order on the current bag.
Replace them by $f+g$ and $f-g$, where $f$ is their average and the
perturbation $g$ vanishes on that bag. The combined cut law on the
bag is unchanged: successive threshold masses are affine in the
ordered heights. This permits different continuations into the
children of a decomposition node while preserving consistency.

Every perturbation has a fixed uniform Lipschitz bound. Nevertheless,
repeated updates might appear to increase the particles' Lipschitz constants
without bound. Section~\ref{sec:particles} prevents this by a
numerical invariant: a quadratic bound on sums of the largest
particle differences controls every individual Lipschitz constant.
Averaging two entries and adding opposite perturbations preserves
this bound, which depends only on the fixed number of particles. Thus every
particle remains uniformly Lipschitz after arbitrarily many updates.

\paragraph{Unchanged particles carry uncertainty along bag paths.}
A \emph{reservation} freezes two particles, preventing their use in
later updates until the reservation is released. Once the maximum
number of simultaneous reservations is fixed, a sufficiently large
particle list always leaves two unreserved particles with the same
chosen increasing order: there are at most $k!$ such orders on a bag.
The functions in a reserved pair remain unchanged, so their two cut
measures form a common submeasure of the laws on every bag
along the reserved segment. At its launch bag $A$, the pair $f\pm g$
agrees because $g|_A=0$. At a remote vertex $z$, the values may
separate by $2|g(z)|$.

The relevant quantity is \emph{uncertainty}: after fixing the binary
assignment on a testing bag, take the smaller of the two masses
corresponding to the possible labels at $z$, and sum over all
assignments. The launch pair has uncertainty at least $2|g(z)|$.
Section~\ref{sec:cuts} proves that a common submeasure transports
this uncertainty to separation of the endpoints of a bag path.
The proof inserts the full testing-bag assignment into a Markov
chain and applies the published symmetric Markov-flow theorem.
The chain has at most $2^k$ states per layer, so its loss depends
only on $k$, independently of the number of bags.

\paragraph{Two reservation schedules.}
Section~\ref{sec:perturbations} first constructs a cut pseudometric
$D_{\mathrm{bag}}$ that separates all pairs in a common bag. For
each bag, only a bounded number of dyadic scales match distances
between its vertices. Particles perturbed using padded partitions
are reserved along the consecutive bags where a scale remains
relevant. This gives both a uniform upper bound on every pair and
a uniform lower bound for pairs sharing a bag.

Section~\ref{sec:anchors} constructs a second cut pseudometric
$D_{\mathrm{anc}}$. Here a reserved pair is attached to a launch
bag, called its \emph{anchor}. A bounded list of anchors is retained
on each descending path of the rooted decomposition. The list
protects entries immediately after certain drops in the distances
from vertices of the current bag to older launch bags. This rule
depends only on the metric and decomposition, before any fresh
perturbation is sampled. The rule protects only a bounded number
of entries, independently of how small the eventual noise parameter
must be.

Bounded bag size thus controls the assignment spaces in transport,
the particle orders, the active scales, and the protected anchor
entries. It also bounds the constants in the padded partitions used
to construct the perturbations.

\paragraph{A younger reservation contradicts the youngest one.}
Section~\ref{sec:lower} combines the two metrics in a compactness
argument. If their sum had no uniform lower bound, the common-bag
comparison would allow us to normalize a sequence of poorly
separated pairs so that one endpoint stays a fixed distance from
the least common ancestor bag. Along the descending path to that
endpoint, the distances from the bags to the endpoint fall to zero.
Passing to a limit along a nonprincipal ultrafilter, only boundedly
many reservations can survive from bags at a definite distance to
bags approaching the endpoint. Choose the youngest such reservation.

Its perturbation contains two independent parts: a random multiple
of distance to its launch bag, and local noise near the endpoint.
Together they force every testing bag still at a definite distance
to contain vertices $p,q$, with $p$ nearest the endpoint, whose
mutual distance lies between fixed multiples of the endpoint's
distances to the testing bag and to the launch bag. Otherwise the
noise produces a positive height gap and hence separation by
uncertainty transport. The proof uses this within-bag distance in
the list rule to protect a younger successor. That successor also
survives toward the endpoint, contradicting the choice of the
youngest reservation.

The resulting positive lower constant depends only on $k$.
Section~\ref{sec:conclusion} adds the two cut pseudometrics and expands their
finite cut representation as in~\eqref{eq:intro-cut-representation},
completing the embedding in Theorem~\ref{thm:main}.
